% Generated draft: requires mathematical and source review.
\documentclass[11pt]{article}
\usepackage[margin=25mm]{geometry}
\usepackage{fontspec}
\setmainfont{Latin Modern Roman}
\usepackage{amsmath,amssymb,mathtools,unicode-math}
\setmathfont{Latin Modern Math}
\usepackage{enumitem,needspace,xurl,hyperref}
\setlength{\parindent}{0pt}
\setlength{\parskip}{5pt}
\setlength{\emergencystretch}{3em}
\newcommand{\N}{\mathbb{N}}
\newcommand{\Z}{\mathbb{Z}}
\newcommand{\Q}{\mathbb{Q}}
\newcommand{\R}{\mathbb{R}}
\newcommand{\C}{\mathbb{C}}
\newcommand{\F}{\mathbb{F}}
\newcommand{\A}{\mathbb{A}}
\newcommand{\PP}{\mathbb{P}}
\newcommand{\E}{\mathbb{E}}
\newcommand{\eps}{\varepsilon}
\newcommand{\dd}{\,\mathrm{d}}
\begin{document}
\section*{UnsolvedMath ID: 30002509}
\textbf{Source problem code:} catalogue entry
\subsection{Definitions and mathematical statement}
A spine of a topological space \(Y\) is understood here as a subspace \(S\subseteq Y\) admitting a strong deformation retraction: a continuous map \(H:Y\times[0,1]\to Y\) such that
\[
H(y,0)=y,\qquad H(y,1)\in S,\qquad H(s,t)=s
\]
for every \(y\in Y\), \(s\in S\), and \(t\in[0,1]\). If a group \(\Gamma\) acts on \(Y\), a \(\Gamma\)-equivariant spine additionally requires \(S\) to be \(\Gamma\)-invariant and
\[
H(\gamma y,t)=\gamma H(y,t)
\quad(\gamma\in\Gamma,\ y\in Y,\ t\in[0,1]).
\]
This is an explicit interpretation of the word ``spine''; the supplied record does not specify its precise deformation-retraction, dimensional, or cocompactness requirements.

For a closed, connected, oriented surface \(\Sigma_g\) of genus \(g\), its Teichm\"uller space \(\mathcal T_g\) is the space of marked conformal structures: pairs \((R,f)\), where \(R\) is a Riemann surface and \(f:\Sigma_g\to R\) is an orientation-preserving marking, modulo conformal isomorphisms compatible with the markings up to isotopy. Under the customary closed-surface interpretation with integer \(g\geq2\), the first requested construction is a spine \(S_g\subseteq\mathcal T_g\) satisfying
\[
\dim S_g=4g-5.
\]
The genus range and absence of punctures or boundary are conventional interpretations, not explicit hypotheses of the supplied statement. The record also does not explicitly require equivariance under the mapping class group \(\operatorname{Mod}(\Sigma_g)\), the group of isotopy classes of orientation-preserving self-homeomorphisms of \(\Sigma_g\). Whether such equivariance and a compact quotient are intended must be clarified; these conditions materially affect the target.

The arithmetic request in the supplied record is: for every nonuniform arithmetic subgroup \(\Gamma<G\) acting on the associated symmetric space \(X\), construct a \(\Gamma\)-equivariant spine \(S\subseteq X\) with
\[
\dim S=\operatorname{vcd}(\Gamma)
       =\dim X-\operatorname{rank}_{\Q}(\Gamma).
\]
Here ``nonuniform'' conventionally means that \(\Gamma\) is a lattice with noncompact quotient, and \(\operatorname{vcd}(\Gamma)\) denotes the cohomological dimension over \(\Z\) of a torsion-free subgroup of finite index, when such a subgroup exists. The notation \(\operatorname{rank}_{\Q}(\Gamma)\) requires a specified ambient algebraic group defined over \(\Q\); in that setting the rational rank is the dimension of a maximal \(\Q\)-split torus of that algebraic group. The record supplies neither this algebraic group nor the structural hypotheses on \(G\) and \(X\) needed to interpret the displayed dimension identity. That identity is therefore retained as part of the source formulation, with its hypotheses unresolved. The required notion of dimension, any cellular structure, and any compact-quotient requirement also remain unspecified.
\subsection{Short English statement}
The supplied problem asks for a spine of Teichm\"uller space of dimension \(4g-5\), and, in an arithmetic setting, an equivariant spine whose dimension equals the virtual cohomological dimension. The intended genus range, mapping class group equivariance, quotient compactness, dimension convention, and arithmetic hypotheses require clarification. This is a statement-only draft requiring human mathematical review.
\subsection{Sources}
\begin{itemize}
\item Supplied local UnsolvedMath registry record, permanent ID \(30002509\), exact source problem code \texttt{catalogue entry}, title ``Optimal-Dimensional Spines for Teichm\"uller and Arithmetic Spaces.'' The record attributes the problem to the 2014 Oberwolfach workshop report; no individual proposer is supplied.
\item Supplied source citation: \emph{New Trends in Teichm\"uller Theory and Mapping Class Groups} (2014), Oberwolfach Reports, pp.~395--452, DOI \texttt{10.4171/owr/2014/07}. \url{https://doi.org/10.4171/owr/2014/07}. The background describes these reports as documenting mathematical workshops and their open-problem discussions; no problem-specific page is supplied.
\item Supplied literature source: ``Thurston spine dimension estimates.'' \url{https://arxiv.org/abs/2310.15618}.
\item Supplied literature source: ``Decorated Teichmuller spines.'' \url{https://arxiv.org/abs/2409.04392}.
\item Literature assessment dated exactly \texttt{2026-08-21}: the supplied curation labels the problem open and reports that spines of the expected virtual-cohomological dimension are known in low codimension and for related decorated spaces, while no general spine of dimension \(4g-5\), or corresponding optimal arithmetic construction, is known. This is supplied dated curation derived from the OWR material, not an independently verified literature assessment; its time-sensitive status and cited claims require independent checking.
\item The registry labels the statement \texttt{corrected\_verified} and says that mathematical content was unchanged while notation, typography, and wording were normalized. These are supplied editorial labels and claims, not independent mathematical verification of this draft.
\end{itemize}
\end{document}
